% =====================================================================
%  CHECKPOINT 02 -- state after block 2 (text, lists and the preamble)
%  Copy this over latex-workshop.tex if you need to catch up.
%
%  House style in these files: one sentence per line, \label on the
%  same line as its \section, four spaces per indent.
% =====================================================================

\documentclass[a4paper,11pt]{article}

% ---------------------------------------------------------------------
%  Preamble
% ---------------------------------------------------------------------
\usepackage[T1]{fontenc}        % proper 8-bit output encoding
\usepackage[english]{babel}     % hyphenation and language conventions
\usepackage{newtxtext}          % Times for the text
\usepackage{parskip}            % blank line between paragraphs, no indent

% TODO 3.1  Add two lines here, in this order:
%
%             \usepackage{amsmath}
%             \usepackage{newtxmath}
%
%           amsmath is where almost everything in block 3 comes from --
%           align, cases, \text and sane spacing.
%           newtxmath makes the MATHS use the same Times design as the
%           text. Try compiling block 3 without it once: Times prose
%           with Computer Modern formulas is a mismatch you cannot
%           unsee afterwards.

\title{How Fast Does Your Coffee Go Cold?}
\author{Your Name}
\date{\today}

\begin{document}

\maketitle
\tableofcontents

% =====================================================================
\section{Introduction}\label{sec:introduction}

A freshly poured coffee is too hot to drink, and a forgotten one is too cold.
Somewhere in between there is a window of a few minutes in which the drink is actually pleasant,
and this report sets out to find it.
The underlying physics is \emph{Newton's law of cooling},
which says that a body loses heat in proportion to how much hotter it is than its surroundings.
Three \textbf{everyday vessels} are compared,
referred to throughout by the short labels \texttt{mug}, \texttt{glass} and \texttt{flask}.

The question is worth asking precisely because the model is so thin.
It has exactly one free parameter per vessel,
so a single afternoon is enough to test it properly.

This report sets out to:

\begin{itemize}
    \item state the cooling model and what its one parameter means,
    \item compare the cooling curves of the three vessels, and
    \item work out how long each of them keeps a drink in the pleasant range.
\end{itemize}

% =====================================================================
\section{Theory}\label{sec:theory}

% TODO 3.2  Rewrite the description list below so that the symbols are
%           typeset as maths: $T_0$, $T_\mathrm{env}$ and $\tau$ rather
%           than upright T_0, T_env and tau. Note what changes about
%           the shapes of the letters, and why that matters.

\begin{description}
    \item[T\_0] temperature of the drink at the start
    \item[T\_env] temperature of the room
    \item[tau] time constant of the vessel: the time it takes to close
        roughly two thirds of the gap to room temperature
\end{description}

% TODO 3.3  Write Newton's law of cooling as a NUMBERED display
%           equation (the equation environment), and give it
%           \label{eq:cooling}:
%
%             T(t) = T_env + (T_0 - T_env) * e^(-t/tau)
%
%           Use \mathrm{} for the "env" subscript and for the e of the
%           exponential, and \tau for the time constant.

% TODO 3.4  Below it, rearrange that equation for the waiting time in
%           an align environment: one line isolating the exponential, a
%           second line with the result
%
%             t = -tau * ln( (T - T_env) / (T_0 - T_env) )
%
%           Align the two lines on their equals signs using &, and give
%           the second line \label{eq:waiting}.

% TODO 3.5  Write a sentence underneath that refers back to both
%           equations by number, using \eqref{eq:cooling} and
%           \eqref{eq:waiting}. Never type the number "1" yourself.

% TODO 3.6  Finally, classify the drink as a case distinction, using the
%           cases environment inside a display equation:
%
%             verdict(T) = too hot     if T > 65 degC
%                          drinkable   if 50 degC <= T <= 65 degC
%                          too cold    if T < 50 degC
%
%           Everything here that is a word -- "too hot", "if" -- is
%           text, not maths: wrap it in \text{}. Use \quad to keep the
%           "if" off the condition. For the temperatures, write
%           $65\,^{\circ}\mathrm{C}$; there is a tidier way, and you
%           will meet it at the end of block 4.

% =====================================================================
\section{Setup}\label{sec:setup}

Each vessel starts with 300 ml of water at 85 degrees Celsius.
A thermometer sits in the centre of the liquid, clear of the walls.
The temperature is taken every 30 seconds for one hour.

The room stays closed and unventilated throughout, at 21 degrees Celsius.

The procedure is as follows:

\begin{enumerate}
    \item Heat 300 ml of water to 85 degrees Celsius.
    \item Pour it into the vessel and start the clock.
    \item Suspend the thermometer in the centre of the liquid.
    \item Record the temperature every 30 seconds for one hour.
\end{enumerate}

% =====================================================================
\section{Results}\label{sec:results}

% =====================================================================
\section{Conclusion}\label{sec:conclusion}

\end{document}

% =====================================================================
%  BLOCK 3 -- MATHS
% =====================================================================
%
%  TODO 3.1 is in the preamble; 3.2 to 3.6 are in the Theory section,
%  in the order you meet them reading downwards.
%
%  THE TWO RULES WORTH REMEMBERING
%    * Anything that is a quantity belongs in maths mode, even in the
%      middle of a sentence: "the time constant $\tau$", not "tau".
%    * Anything that is a word belongs in \text{} or \mathrm{}, even in
%      the middle of a formula: $T_\mathrm{env}$, not $T_{env}$.
%
%  IF YOU FINISH EARLY -- typesetting drills
%  These do not belong in the report; put them in a scratch file, or at
%  the end of the document and delete them afterwards.
%
%    (a) The logistic map          x_{t+1} = k x_t (1 - x_t)
%
%    (b) Shannon entropy           H = - sum from i=0 to N of
%                                      p(x_i) log p(x_i)
%
%    (c) Hard mode -- the inverse Laplace transform
%                                  L^-1{F} = lim as T -> infinity of
%                                      1/(2 pi i) * integral from
%                                      gamma - iT to gamma + iT of
%                                      e^(st) F(s) ds
%
%  Symbol you do not know? https://detexify.kirelabs.org/ -- draw it,
%  and it tells you the command.
% =====================================================================
